On ne considere dans cette section que le probleme local de l'evolution
de la distribution des electrons dans l'espace des impulsions. Sans
le mentionner explicitement pour alleger les notations, les valeurs
des fonctions et des flux sont donc evalues a la meme position radiale,
soit pour $f$ la demi-grille radiale $r_{l+1/2}$ et polo\"{\i}dale
$\theta_{m+1/2}$ et pour les flux la grille entiere $r_{l}$ et polo\"{\i}dale
$\theta_{m}$ comme indique dans la section dediee a la definition
des differents maillages. On obtient donc \[
\frac{\partial\left(rp^{2}f_{0}\right)}{\partial t}+r\frac{\partial\left(p^{2}S_{p}\right)}{\partial p}-rp\frac{\partial\left(\sqrt{1-\xi^{2}}S_{\xi}\right)}{\partial\xi}+p^{2}\frac{\partial\left(rS_{r}\right)}{\partial r}+p^{2}\frac{\partial S_{\theta}}{\partial\theta}=0\]
\begin{eqnarray*}
 &  & r_{l+1/2}p_{i+1/2}^{2}\frac{f_{0,i+1/2,j+1/2}^{(k+1)}-f_{0,i+1/2,j+1/2}^{(k)}}{\Delta t}+r_{l+1/2}\left.\frac{\partial\left(p^{2}S_{p}\right)}{\partial p}\right|_{i+1/2,j+1/2}^{(k+1)}\\
 &  & -r_{l+1/2}p_{i+1/2}\left.\frac{\partial\left(\sqrt{1-\xi^{2}}S_{\xi}\right)}{\partial\xi}\right|_{i+1/2,j+1/2}^{(k+1)}=0\end{eqnarray*}
 avec \begin{eqnarray*}
S_{p} & = & -D_{pp}\frac{\partial f_{0}}{\partial p}+D_{p\xi}\frac{\sqrt{1-\xi^{2}}}{p}\frac{\partial f_{0}}{\partial\xi}+F_{p}f_{0}\\
S_{\xi} & = & -D_{\xi p}\frac{\partial f_{0}}{\partial p}+D_{\xi\xi}\frac{\sqrt{1-\xi^{2}}}{p}\frac{\partial f_{0}}{\partial\xi}+F_{\xi}f_{0}\end{eqnarray*}
 Sachant que les termes $D_{p\xi}$ et $D_{\xi p}$ ne concernent
que les perturbations a la solution Maxwellienne, et que dans la region
concernee de l'espace des impulsions ou celles-ci interviennent la
distribution peut etre de forme quelconque, il est utile dans ce cas
d'effectuer un traitement particulier, en centrant simplement les
differences. On verra en effet plus loin qu'il n'est pas utile d'effectuer
une interpolation specifique entre grille et demi-grille pour les
derivees croisees afin d'assurer une precision numerique correcte.
Ainsi, \begin{eqnarray*}
\frac{\partial\left(p^{2}S_{p}\right)}{\partial p} & = & \frac{\partial}{\partial p}\left(-p^{2}D_{pp}\frac{\partial f_{0}}{\partial p}+p^{2}F_{p}f_{0}\right)+\sqrt{1-\xi^{2}}\frac{\partial}{\partial p}\left(pD_{p\xi}\right)\frac{\partial f_{0}}{\partial\xi}+\sqrt{1-\xi^{2}}pD_{p\xi}\frac{\partial^{2}f_{0}}{\partial p\partial\xi}\\
\frac{\partial\left(\sqrt{1-\xi^{2}}S_{\xi}\right)}{\partial\xi} & = & \frac{\partial}{\partial\xi}\left(D_{\xi\xi}\frac{1-\xi^{2}}{p}\frac{\partial f_{0}}{\partial\xi}+\sqrt{1-\xi^{2}}F_{\xi}f_{0}\right)-\frac{\partial}{\partial\xi}\left(\sqrt{1-\xi^{2}}D_{\xi p}\right)\frac{\partial f_{0}}{\partial p}-\sqrt{1-\xi^{2}}D_{\xi p}\frac{\partial^{2}f_{0}}{\partial\xi\partial p}\end{eqnarray*}
 et l'on obtient finalement \begin{eqnarray*}
 &  & r_{l+1/2}p_{i+1/2}^{2}\frac{f_{0,i+1/2,j+1/2}^{(k+1)}-f_{0,i+1/2,j+1/2}^{(k)}}{\Delta t}\\
 &  & +r_{l+1/2}\frac{-p_{i+1}^{2}D_{pp,i+1,j+1/2}\left.\frac{\partial f_{0}}{\partial p}\right|_{i+1,j+1/2}^{(k+1)}+p_{i+1}^{2}F_{p,i+1,j+1/2}f_{0,i+1,j+1/2}^{(k+1)}}{\Delta p_{i+1/2}}\\
 &  & +r_{l+1/2}\frac{p_{i}^{2}D_{pp,i,j+1/2}\left.\frac{\partial f_{0}}{\partial p}\right|_{i,j+1/2}^{(k+1)}-p_{i}^{2}F_{p,i,j+1/2}f_{0,i,j+1/2}^{(k+1)}}{\Delta p_{i+1/2}}\\
 &  & +r_{l+1/2}\sqrt{1-\xi_{j+1/2}^{2}}p_{i+1/2}D_{p\xi,i+1/2,j+1/2}\left.\frac{\partial^{2}f_{0}}{\partial p\partial\xi}\right|_{i+1/2,j+1/2}^{(k+1)}\\
 &  & +r_{l+1/2}\sqrt{1-\xi_{j+1/2}^{2}}\left.\frac{\partial\left(pD_{p\xi}\right)}{\partial p}\right|_{i+1/2,j+1/2}\left.\frac{\partial f_{0}}{\partial\xi}\right|_{i+1/2,j+1/2}^{(k+1)}\\
 &  & -r_{l+1/2}\frac{D_{\xi\xi,i+1/2,j+1}\left(1-\xi_{j+1}^{2}\right)\left.\frac{\partial f_{0}}{\partial\xi}\right|_{i+1/2,j+1}^{(k+1)}+p_{i+1/2}\sqrt{1-\xi_{j+1}^{2}}F_{\xi,i+1/2,j+1}f_{0,i+1/2,j+1}^{(k+1)}}{\Delta\xi_{j+1/2}}\\
 &  & +r_{l+1/2}\frac{D_{\xi\xi,i+1/2,j}\left(1-\xi_{j}^{2}\right)\left.\frac{\partial f_{0}}{\partial\xi}\right|_{i+1/2,j}^{(k+1)}+p_{i+1/2}\sqrt{1-\xi_{j}^{2}}F_{\xi,i+1/2,j}f_{0,i+1/2,j}^{(k+1)}}{\Delta\xi_{j+1/2}}\\
 &  & +r_{l+1/2}p_{i+1/2}\sqrt{1-\xi_{j+1/2}^{2}}D_{\xi p,i+1/2,j+1/2}\left.\frac{\partial^{2}f_{0}}{\partial\xi\partial p}\right|_{i+1/2,j+1/2}^{(k+1)}\\
 &  & +r_{l+1/2}p_{i+1/2}\left.\frac{\partial}{\partial\xi}\left(\sqrt{1-\xi^{2}}D_{\xi p}\right)\right|_{i+1/2,j+1/2}\left.\frac{\partial f_{0}}{\partial p}\right|_{i+1/2,j+1/2}^{(k+1)}\\
 & = & 0\end{eqnarray*}


ou l'indice $(k)$ indique que la variable est prise au temps $t=k\Delta t,$
$\Delta t$ etant le pas de temps. Le terme de droite est pris au
temps $(k+1)$, ce qui correspond a un schema totalement implicite
ou temps inverse, inconditionnellement stable numeriquement, point
qui sera rediscute dans un chapitre suivant.

Les expressions discretes des derivees partielles prennent la forme
suivante, \[
\left\{ \begin{array}{l}
\left.\frac{\partial f_{0}}{\partial p}\right|_{i+1,j+1/2}^{(k+1)}=\frac{f_{0,i+3/2,j+1/2}^{(k+1)}-f_{0,i+1/2,j+1/2}^{(k+1)}}{\Delta p_{i+1}}\\
\left.\frac{\partial f_{0}}{\partial p}\right|_{i+1/2,j+1/2}^{(k+1)}=\frac{f_{0,i+3/2,j+1/2}^{(k+1)}-f_{0,i-1/2,j+1/2}^{(k+1)}}{\Delta p_{i+1}+\Delta p_{i}}\\
\left.\frac{\partial f_{0}}{\partial p}\right|_{i,j+1/2}^{(k+1)}=\frac{f_{0,i+1/2,j+1/2}^{(k+1)}-f_{0,i-1/2,j+1/2}^{(k+1)}}{\Delta p_{i}}\\
\left.\frac{\partial f_{0}}{\partial\xi}\right|_{i+1/2,j+1}^{(k+1)}=\frac{f_{0,i+1/2,j+3/2}^{(k+1)}-f_{0,i+1/2,j+1/2}^{(k+1)}}{\Delta\xi_{j+1}}\\
\left.\frac{\partial f_{0}}{\partial\xi}\right|_{i+1/2,j+1/2}^{(k+1)}=\frac{f_{0,i+1/2,j+3/2}^{(k+1)}-f_{0,i+1/2,j-1/2}^{(k+1)}}{\Delta\xi_{j+1}+\Delta\xi_{j}}\\
\left.\frac{\partial f_{0}}{\partial\xi}\right|_{i+1/2,j}^{(k+1)}=\frac{f_{0,i+1/2,j+1/2}^{(k+1)}-f_{0,i+1/2,j-1/2}^{(k+1)}}{\Delta\xi_{j}}\end{array}\right.\]
 et \begin{eqnarray*}
\left.\frac{\partial^{2}f_{0}}{\partial p\partial\xi}\right|_{i+1/2,j+1/2}^{(k+1)} & = & \left.\frac{\partial^{2}f_{0}}{\partial\xi\partial p}\right|_{i+1/2,j+1/2}^{(k+1)}\\
 & = & \frac{f_{0,i+3/2,j+3/2}^{(k+1)}-f_{0,i+3/2,j-1/2}^{(k+1)}-f_{0,i-1/2,j+3/2}^{(k+1)}+f_{0,i-1/2,j-1/2}^{(k+1)}}{\left(\Delta p_{i+1}+\Delta p_{i}\right)\left(\Delta\xi_{j+1}+\Delta\xi_{j}\right)}\end{eqnarray*}
 tandis que \begin{eqnarray*}
\left.\frac{\partial\left(pD_{p\xi}\right)}{\partial p}\right|_{i+1/2,j+1/2} & = & \frac{p_{i+1}D_{p\xi,i+1,j+1/2}-p_{i}D_{p\xi,i,j+1/2}}{\Delta p_{i+1/2}}\\
\left.\frac{\partial}{\partial\xi}\left(\sqrt{1-\xi^{2}}D_{\xi p}\right)\right|_{i+1/2,j+1/2} & = & \frac{\sqrt{1-\xi_{j+1}^{2}}D_{\xi p,i+1/2,j+1}-\sqrt{1-\xi_{j}^{2}}D_{\xi p,i+1/2,j}}{\Delta\xi_{j+1/2}}\end{eqnarray*}


A noter que les derivees croisees sont symetriques par construction.
Par ailleurs la derivee en $\xi$ de $\sqrt{1-\xi^{2}}D_{\xi p}$
est prise sur la grille des flux pour assurer naturellement la condition
aux limites en $\left|\xi\right|=1.$ Ce point est essentiel pour
l'onde hybride pour laquelle cette region de l'espace des impulsions
joue un role primordial pour la valeur du courant.

Comme la fonction est definie sur la demi-grille et les flux sur la
grille entiere, il faut effectuer une interpolation, puisque certaines
valeurs de $f$ apparaissent sur la grille entiere dans les relations
precedentes. De maniere generale, sans se preoccuper de l'expression
precise du facteur de ponderation $\delta$ qui sera discutee plus
loin, on peut ecrire pour les termes proportionnnels a $D_{pp}$ et
$F_{p}$\[
\left\{ \begin{array}{l}
f_{0,i+1,j+1/2}^{(k+1)}=\left(1-\delta_{p,i+1,j+1/2}\right)f_{0,i+3/2,j+1/2}^{(k+1)}+\delta_{p,i+1,j+1/2}f_{0,i+1/2,j+1/2}^{(k+1)}\\
f_{0,i,j+1/2}^{(k+1)}=\left(1-\delta_{p,i,j+1/2}\right)f_{0,i+1/2,j+1/2}^{(k+1)}+\delta_{p,i,j+1/2}f_{0,i-1/2,j+1/2}^{(k+1)}\end{array}\right.\]
 et ceux proportionnels a $D_{\xi\xi}$ et $F_{\xi}$

\[
\left\{ \begin{array}{l}
f_{0,i+1/2,j+1}^{(k+1)}=\left(1-\delta_{\xi,i+1/2,j+1}\right)f_{0,i+1/2,j+3/2}^{(k+1)}+\delta_{\xi,i+1/2,j+1}f_{0,i+1/2,j+1/2}^{(k+1)}\\
f_{0,i+1/2,j}^{(k+1)}=\left(1-\delta_{\xi,i+1/2,j}\right)f_{0,i+1/2,j+1/2}^{(k+1)}+\delta_{\xi,i+1/2,j}f_{0,i+1/2,j-1/2}^{(k+1)}\end{array}\right.\]


Il est interessant a ce stade des calculs d'exprimer les coefficients
de la matrice qui correspondent a l'equation de Fokker-Plank discretisee
\[
r_{l+1/2}p_{i+1/2}^{2}\frac{f_{0,i+1/2,j+1/2}^{\left(k+1\right)}-f_{0,i+1/2,j+1/2}^{\left(k\right)}}{\Delta t}+r_{l+1/2}\sum_{i^{\prime}=i-1}^{i^{\prime}=i+1}\sum_{j^{\prime}=j-1}^{j^{\prime}=j+1}\overline{\overline{M}}_{p,i^{\prime}+1/2,j^{\prime}+1/2}f_{0,i^{\prime}+1/2,j^{\prime}+1/2}^{(k+1)}=0\]
 On obtient \begin{eqnarray*}
\overline{\overline{M}}_{p,i+3/2,j+3/2} & = & \frac{p_{i+1/2}\sqrt{1-\xi_{j+1/2}^{2}}}{\Delta\xi_{j+1}+\Delta\xi_{j}}\frac{D_{p\xi,i+1/2,j+1/2}}{\Delta p_{i+1}+\Delta p_{i}}\\
 &  & +\frac{p_{i+1/2}\sqrt{1-\xi_{j+1/2}^{2}}}{\Delta\xi_{j+1}+\Delta\xi_{j}}\frac{D_{\xi p,i+1/2,j+1/2}}{\Delta p_{i+1}+\Delta p_{i}}\end{eqnarray*}
\begin{eqnarray*}
\overline{\overline{M}}_{p,i+1/2,j+3/2} & = & \frac{p_{i+1}}{\Delta p_{i+1/2}}\frac{\sqrt{1-\xi_{j+1/2}^{2}}}{\Delta\xi_{j+1}+\Delta\xi_{j}}D_{p\xi,i+1,j+1/2}\\
 &  & -\frac{p_{i}}{\Delta p_{i+1/2}}\frac{\sqrt{1-\xi_{j+1/2}^{2}}}{\Delta\xi_{j+1}+\Delta\xi_{j}}D_{p\xi,i,j+1/2}\\
 &  & -\frac{1-\xi_{j+1}^{2}}{\Delta\xi_{j+1/2}\Delta\xi_{j+1}}D_{\xi\xi,i+1/2,j+1}\\
 &  & -p_{i+1/2}\frac{\sqrt{1-\xi_{j+1}^{2}}}{\Delta\xi_{j+1/2}}\left(1-\delta_{\xi,i+1/2,j+1}\right)F_{\xi,i+1/2,j+1}\end{eqnarray*}
\begin{eqnarray*}
\overline{\overline{M}}_{p,i-1/2,j+3/2} & = & -p_{i+1/2}\frac{\sqrt{1-\xi_{j+1/2}^{2}}}{\Delta\xi_{j+1}+\Delta\xi_{j}}\frac{D_{p\xi,i+1/2,j+1/2}}{\Delta p_{i+1}+\Delta p_{i}}\\
 &  & -p_{i+1/2}\frac{\sqrt{1-\xi_{j+1/2}^{2}}}{\Delta\xi_{j+1}+\Delta\xi_{j}}\frac{D_{\xi p,i+1/2,j+1/2}}{\Delta p_{i+1}+\Delta p_{i}}\end{eqnarray*}
\begin{eqnarray*}
\overline{\overline{M}}_{p,i+3/2,j+1/2} & = & -p_{i+1}^{2}\frac{D_{pp,i+1,j+1/2}}{\Delta p_{i+1}\Delta p_{i+1/2}}\\
 &  & +p_{i+1/2}\frac{\sqrt{1-\xi_{j+1}^{2}}}{\Delta\xi_{j+1/2}}\frac{D_{\xi p,i+1/2,j+1}}{\Delta p_{i+1}+\Delta p_{i}}\\
 &  & -p_{i+1/2}\frac{\sqrt{1-\xi_{j}^{2}}}{\Delta\xi_{j+1/2}}\frac{D_{\xi p,i+1/2,j}}{\Delta p_{i+1}+\Delta p_{i}}\\
 &  & +\frac{p_{i+1}^{2}}{\Delta p_{i+1/2}}\left(1-\delta_{p,i+1,j+1/2}\right)F_{p,i+1,j+1/2}\end{eqnarray*}
\begin{eqnarray*}
\overline{\overline{M}}_{p,i+1/2,j+1/2} & = & \frac{p_{i+1}^{2}}{\Delta p_{i+1/2}\Delta p_{i+1}}D_{pp,i+1,j+1/2}\\
 &  & +\frac{p_{i+1}^{2}}{\Delta p_{i+1/2}}F_{p,i+1,j+1/2}\delta_{p,i+1,j+1/2}\\
 &  & +\frac{p_{i}^{2}}{\Delta p_{i+1/2}\Delta p_{i}}D_{pp,i,j+1/2}\\
 &  & -\frac{p_{i}^{2}}{\Delta p_{i+1/2}}\left(1-\delta_{p,i,j+1/2}\right)F_{p,i,j+1/2}\\
 &  & +\frac{1-\xi_{j+1}^{2}}{\Delta\xi_{j+1/2}\Delta\xi_{j+1}}D_{\xi\xi,i+1/2,j+1}\\
 &  & -p_{i+1/2}\frac{\sqrt{1-\xi_{j+1}^{2}}}{\Delta\xi_{j+1/2}}\delta_{\xi,i+1/2,j+1}F_{\xi,i+1/2,j+1}\\
 &  & +\frac{1-\xi_{j}^{2}}{\Delta\xi_{j}\Delta\xi_{j+1/2}}D_{\xi\xi,i+1/2,j}\\
 &  & +p_{i+1/2}\frac{\sqrt{1-\xi_{j}^{2}}}{\Delta\xi_{j+1/2}}\left(1-\delta_{\xi,i+1/2,j}\right)F_{\xi,i+1/2,j}\end{eqnarray*}
\begin{eqnarray*}
\overline{\overline{M}}_{p,i-1/2,j+1/2} & = & -\frac{p_{i}^{2}}{\Delta p_{i+1/2}\Delta p_{i}}D_{pp,i,j+1/2}\\
 &  & -p_{i+1/2}\frac{\sqrt{1-\xi_{j+1}^{2}}}{\Delta\xi_{j+1/2}}\frac{D_{\xi p,i+1/2,j+1}}{\Delta p_{i+1}+\Delta p_{i}}\\
 &  & +p_{i+1/2}\frac{\sqrt{1-\xi_{j}^{2}}}{\Delta\xi_{j+1/2}}\frac{D_{\xi p,i+1/2,j}}{\Delta p_{i+1}+\Delta p_{i}}\\
 &  & -\frac{p_{i}^{2}}{\Delta p_{i+1/2}}\delta_{p,i,j+1/2}F_{p,i,j+1/2}\end{eqnarray*}
\begin{eqnarray*}
\overline{\overline{M}}_{p,i+3/2,j-1/2} & = & -p_{i+1/2}\frac{\sqrt{1-\xi_{j+1/2}^{2}}}{\Delta\xi_{j+1}+\Delta\xi_{j}}\frac{D_{p\xi,i+1/2,j+1/2}}{\Delta p_{i+1}+\Delta p_{i}}\\
 &  & -p_{i+1/2}\frac{\sqrt{1-\xi_{j+1/2}^{2}}}{\Delta\xi_{j+1}+\Delta\xi_{j}}\frac{D_{\xi p,i+1/2,j+1/2}}{\Delta p_{i+1}+\Delta p_{i}}\end{eqnarray*}
\begin{eqnarray*}
\overline{\overline{M}}_{p,i+1/2,j-1/2} & = & -\frac{p_{i+1}}{\Delta p_{i+1/2}}\frac{\sqrt{1-\xi_{j+1/2}^{2}}}{\Delta\xi_{j+1}+\Delta\xi_{j}}D_{p\xi,i+1,j+1/2}\\
 &  & +\frac{p_{i}}{\Delta p_{i+1/2}}\frac{\sqrt{1-\xi_{j+1/2}^{2}}}{\Delta\xi_{j+1}+\Delta\xi_{j}}D_{p\xi,i,j+1/2}\\
 &  & -\frac{1-\xi_{j}^{2}}{\Delta\xi_{j+1/2}\Delta\xi_{j}}D_{\xi\xi,i+1/2,j}\\
 &  & +p_{i+1/2}\frac{\sqrt{1-\xi_{j}^{2}}}{\Delta\xi_{j+1/2}}\delta_{\xi,i+1/2,j}F_{\xi,i+1/2,j}\end{eqnarray*}
\begin{eqnarray*}
\overline{\overline{M}}_{p,i-1/2,j-1/2} & = & p_{i+1/2}\frac{\sqrt{1-\xi_{j+1/2}^{2}}}{\Delta\xi_{j+1}+\Delta\xi_{j}}\frac{D_{p\xi,i+1/2,j+1/2}}{\Delta p_{i+1}+\Delta p_{i}}\\
 &  & +p_{i+1/2}\frac{\sqrt{1-\xi_{j+1/2}^{2}}}{\Delta\xi_{j+1}+\Delta\xi_{j}}\frac{D_{\xi p,i+1/2,j+1/2}}{\Delta p_{i+1}+\Delta p_{i}}\end{eqnarray*}


On voit apparaitre dans chaque coefficient certains termes identiques
qu'il faut evaluer pour chaque mecanisme physique en jeu (collisions,
champ electrique constant, ondes RF, $\cdots$): \[
\overline{\overline{D}}_{p}\rightarrow\left\{ \begin{array}{l}
D_{pp,i+1,j+1/2}\\
D_{pp,i,j+1/2}\\
D_{p\xi,i+1,j+1/2}\\
D_{\xi p,i+1/2,j+1}\\
D_{p\xi,i+1/2,j+1/2}\\
D_{\xi p,i+1/2,j+1/2}\\
D_{p\xi,i,j+1/2}\\
D_{\xi p,i+1/2,j}\\
D_{\xi\xi,i+1/2,j+1}\\
D_{\xi\xi,i+1/2,j}\end{array}\right.\]
 et \[
\overline{\overline{F}}_{p}\rightarrow\left\{ \begin{array}{l}
F_{p,i+1,j+1/2}\\
F_{p,i,j+1/2}\\
F_{\xi,i+1/2,j+1}\\
F_{\xi,i+1/2,j}\end{array}\right.\]

